\section{Scope, Chronology, and Evidence}
\label{sec:tertullian-apparatus}

\subsection{Identity and bounds}

\textbf{Name.} This study uses \emph{Tertullian}, Latin \emph{Tertullianus},
with the conventional geographical description ``of Carthage.'' The
third-person closing sentence of \emph{On the Veiling of Virgins} 17 names
``Septimius Tertullianus'' and says the little work is his. This is a
transmitted internal attribution, not a modern signature. The expanded
``Quintus Septimius Florens Tertullianus'' rests on variable medieval paratext
and is not treated as a complete contemporary legal name. ``Saint'' is not
used: no competent canonization, stable early cult, or received feast was
located. ``Father'' may describe historical influence but is not made an
ecclesial status claim.

\textbf{Life range.} Birth and death are unknown. Current scholarly estimates
place birth variously around 150/160 or as late as about 170. By 197 he had
become an educated Christian author. \emph{To His Wife} securely establishes
that he was married by that work's uncertain late-190s or early-200s
composition; it does not date the wedding. \emph{To Scapula} was written after
the eclipse of 14 August 212; other works may follow, but none supplies a
secure final year. Death is conventionally placed about 220 or 225, largely by
fitting relative literary chronology to Jerome's late report of extreme old
age. This study therefore treats \emph{fl.~c.~196/197--212 or later} as the
controlled range and leaves both endpoints of the life open. Dates use
proleptic Common Era.

\textbf{Places and jurisdiction.} Carthage in Roman Africa Proconsularis is the
secure center of authorship and ecclesial life. Rome is central to several
arguments and later traditions but is not a securely documented residence.
Phrygia and Asia belong to the history of the New Prophecy; no journey there is
claimed. Historical place names do not establish modern ethnicity or political
identity. The civil setting is the Roman Empire under the Severan dynasty and
its provincial administration; the ecclesial setting is the Christian
community of Carthage in relation to African, Roman, Asian, and other churches.

\textbf{Language and state of life.} Tertullian's surviving corpus is Latin; he
reports Greek works now lost, knew Greek Scripture and learning, and is known to
Eusebius through a Greek translation of the \emph{Apologeticum}. His biblical
Latin may preserve Old Latin forms or his own renderings. He was a convert from
a pagan past and a married Christian. Ordination as presbyter is Jerome's late
report and remains doubtful beside Tertullian's lay self-placement. The corpus
shows explicit adherence to the New Prophecy and a breach with Christians he
calls \emph{psychici}; it does not preserve a dated institutional transfer or
named excommunication.

\subsection{Reader, question, and thesis boundary}

The intended reader is a serious general reader, cleric, student, or researcher
who needs one sustained life without a fabricated chronology. The governing
question is how a married Christian intellectual in Carthage made Latin prose
carry an integrated theology of Church, body, truth, and discipline, and why
his confidence that the Paraclete demanded increasing rigor led toward
isolation.

The historical thesis is that Tertullian's achievement and estrangement grew
from related habits: material and linguistic exactness, refusal of divided
allegiance, adversarial boundary-making, and insistence that belief take bodily
form. The publication does not diagnose his personality, settle his eternal
state, decide every work date, reconstruct a lost Marcionite Bible, issue an
ecclesial judgment, or turn later Catholic citation into approval of the whole
corpus.

Included are the controlled life, Carthaginian setting, principal works by
genre and development, apologetics, worship, Scripture and rule of faith,
Trinitarian and Christological language, anthropology and resurrection,
marriage and gender, New Prophecy allegiance, ministry and absolution dispute,
manuscript survival, spuria and lost works, ancient and Catholic reception, and
memorable traditions. Excluded are a line-by-line commentary on every treatise,
complete critical-edition collation, full Vetus Latina citation inventory,
exhaustive social histories of African Christianity, Marcion, or the New
Prophecy, and unsupported cult or relic claims.

\subsection{Evidence key}

\begin{evidencelegend}
\evidencelegendrow{A}{Subject-authored evidence in Tertullian's works, controlled
by genre, audience, revision, relative date, and manuscript transmission.
Self-reference remains literary self-presentation rather than a neutral file.}
\evidencelegendrow{N}{Contemporary or near-contemporary external witness, such
as the African martyr records or early literary reception. Proximity does not
make one author eyewitness to Tertullian.}
\evidencelegendrow{E}{Early received tradition, especially Eusebius, Lactantius,
Jerome, Augustine, and Vincent; it establishes early memory and sometimes
preserves lost reports without becoming contemporary biography.}
\evidencelegendrow{L}{Later manuscript, devotional, visual, local, or popular
tradition, including expanded names, quotations, and inferred occupations. It
is evidence first for its transmitting period.}
\evidencelegendrow{M}{Official ecclesial reception, limited to the act's object:
for example a Catechism citation or papal catechesis, neither of which
canonizes the author or approves every work.}
\evidencelegendrow{K}{Named modern critical reconstruction, with premises and
alternatives exposed; this includes chronology, ethnicity, institutional
status, opponent identification, attribution, and manuscript relationships.}
\evidencelegendrow{S}{Source-grounded synthesis by this publication, never
attributed to an ancient witness.}
\end{evidencelegend}

\subsection{Detailed chronology}

\begin{biotimeline}
\biotimelinerow{17 July 180}{Carthage}{Scillitan Christians are tried by the
proconsul and executed. This anchors a Latin African Christian setting before
Tertullian's corpus, not his presence.}{N; contextual event}
\biotimelinerow{c.~155--170}{Unknown; Africa often inferred}{Possible birth
range derived from competing modern career chronologies. No ancient birth date
or secure birthplace survives.}{K; broad and disputed}
\biotimelinerow{before 197}{Probably Carthage}{Literary education and Christian
adherence precede the surviving public apologetic. The human agents, place,
date, and sequence of conversion and baptism remain unknown.}{A/K; states
strong, dates open}
\biotimelinerow{19 February 197 and after}{Carthage}{\emph{To the Nations}
alludes to the defeat of Clodius Albinus; the \emph{Apologeticum} substantially
reworks the case and is commonly dated later in 197.}{A/K; unusually strong
terminus}
\biotimelinerow{late 190s--early 200s}{Carthage}{Early apologetic, moral, and
catechetical works usually include \emph{To the Martyrs}, \emph{On Spectacles},
\emph{On Prayer}, \emph{On Baptism}, \emph{On Repentance}, and \emph{To His
Wife}; their internal order varies by reconstruction. \emph{To His Wife}
securely establishes marriage by its composition, not a wedding date.}{A/K;
family and sequence approximate}
\biotimelinerow{202/203}{Carthage}{Perpetua, Felicity, and companions are
imprisoned and martyred. Tertullian's editorship of their \emph{Passion} is a
modern hypothesis.}{N; context secure, attribution K}
\biotimelinerow{late second--early third century}{Asia Minor / later
transmission}{The hostile New Prophecy materials excerpted in Eusebius
\emph{Ecclesiastical History} 5.16--19 arise nearer the Asian controversy.
They do not narrate Tertullian's African adherence.}{E; early fragments in a
hostile, mediated dossier}
\biotimelinerow{early 200s}{Carthage}{Anti-heretical work expands:
\emph{Prescription against Heretics}, works against Hermogenes and the
Valentinians, and successive forms of \emph{Against Marcion}. Exact dating
depends on cross-reference and doctrinal sequence.}{A/K; relative}
\biotimelinerow{first decade of third century}{Carthage}{\emph{Against Marcion}
passes through the author-described hurried first work, illicitly circulated
second form, and expanded extant recension.}{A; stages secure, dates K}
\biotimelinerow{c.~207--212 or later}{Carthage}{Works increasingly invoke the
Paraclete and New Prophecy in disputes over veiling, fasting, marriage,
forgiveness, Christology, and ecclesial authority. No single conversion year
controls the transition.}{A/K; allegiance secure, chronology disputed}
\biotimelinerow{c.~208--213 or later}{Carthage}{Major theological works commonly
placed in this mature period include \emph{On the Soul}, \emph{On the Flesh of
Christ}, \emph{On the Resurrection of the Flesh}, and \emph{Against Praxeas}.
Individual ranges and order remain contested.}{A/K; broad placement}
\biotimelinerow{14 August 212 and after}{Carthage}{The eclipse recalled in
\emph{To Scapula} controls its date; Tertullian warns the proconsul against
coercion and persecution.}{A/K; strong terminus}
\biotimelinerow{after 212}{Carthage presumed}{\emph{On Monogamy}, \emph{On
Fasting}, and \emph{On Modesty} are ordinarily among the late rigorist works.
Their dates and the identity of the bishop mocked in \emph{On Modesty} remain
open.}{A/K; relative}
\biotimelinerow{after 212; often c.~220--225}{Unknown}{Death. Jerome says only
that Tertullian lived to advanced old age. No controlled cause, place, final
communion, burial, or cult record survives.}{E/K; endpoint unknown}
\biotimelinerow{c.~249--258}{Carthage}{Cyprian's writings receive Tertullian's
language and themes while placing rigor within Catholic episcopal communion.}{N/E;
literary dependence strong}
\biotimelinerow{early fourth century}{Roman Empire}{Lactantius praises
Tertullian's learning and apologetic completeness while criticizing his
obscurity and polish.}{E; direct literary reception}
\biotimelinerow{early fourth century}{Caesarea}{Eusebius transmits a Greek
translation of the \emph{Apologeticum} and calls Tertullian learned in Roman
law.}{E; reputation, not career record}
\biotimelinerow{c.~392/393}{Bethlehem}{Jerome's \emph{On Illustrious Men} 53
records presbyterate, centurion father, Roman envy, Montanist breach, Cyprianic
anecdote, and old age.}{E; indispensable late biography}
\biotimelinerow{420s}{North Africa}{Augustine reports surviving Carthaginian
Tertullianists, their dwindling return to Catholic communion, and transfer of a
basilica.}{E; fifth-century group}
\biotimelinerow{c.~434}{Gaul}{Vincent of L\'erins praises Tertullian's Latin
genius and makes his later error a warning against private novelty.}{E; received
evaluation}
\biotimelinerow{late eighth century}{Frankish realm / Paris witness}{Earliest
located surviving manuscript item: fragment of \emph{Against the Jews}, now
BnF lat. 13047.}{L/K; material transmission}
\biotimelinerow{ninth--twelfth centuries}{Western Europe}{BnF lat. 1622, the
Codex Agobardinus, and Troyes 523 preserve distinct corpus routes. Troyes 523
is the sole extant representative of the Trecense family and carries five
Tertullian works; that does not make it the sole witness to each work.}{L/K;
manuscript evidence}
\biotimelinerow{sixteenth century}{Basel, Paris, and other print centers}{Early
printed editions gather the corpus and preserve readings from Corbie manuscripts
subsequently lost.}{L/K; edition history}
\biotimelinerow{1954}{Turnhout}{\emph{Corpus Christianorum, Series Latina}
volumes 1--2 publish the principal modern collected critical corpus with
work-specific editors.}{K; critical-edition milestone}
\biotimelinerow{30 May 2007}{Vatican City}{Benedict XVI praises Tertullian's
Latin theological contribution and identifies rigorist isolation as the drama
of his life.}{M; papal catechesis}
\biotimelinerow{2026}{International scholarship}{Binder's study of Jews in
Tertullian's literary imagination and other current work continue revision of
identity, rhetoric, chronology, and reception.}{K; mutable scholarship}
\end{biotimeline}

\subsection{Chronology controls and competing reconstructions}

The chronology does not derive from a dated ancient life. It combines three
harder anchors---the Albinus allusion after 19 February 197, the eclipse after
14 August 212, and authorial cross-references---with relative development,
historical names, adversarial targets, and style. The order of the three forms
of \emph{Against Marcion} is author-attested; their exact years are not. The
place of \emph{Against Praxeas}, \emph{On the Crown}, \emph{On Flight},
\emph{Scorpiace}, \emph{On the Resurrection of the Flesh}, and the rigorist
works varies among modern tables.

Timothy Barnes's historical revision remains foundational but has itself been
revised. Older tables often begin life around 155/160 and extend writing or life
into the 220s. Jared Secord argues for a birth perhaps around 170 and a compact
literary career roughly 196--212. Current reference works continue to print
broader ranges. This study does not force apparent theological change to yield
one exact calendar. Local passages retain their event dates; the appendix owns
the disputed whole-life ranges.

\subsection{Corpus and method}

The principal subject-authored corpus inspected includes \emph{To the Nations},
\emph{Apologeticum}, \emph{To Scapula}, \emph{To the Martyrs}, \emph{On
Spectacles}, \emph{On Idolatry}, \emph{On the Crown}, \emph{Scorpiace},
\emph{On Flight in Persecution}, \emph{On Baptism}, \emph{On Prayer},
\emph{On Repentance}, \emph{On Patience}, \emph{To His Wife}, \emph{On the
Apparel of Women}, \emph{On the Veiling of Virgins}, \emph{Exhortation to
Chastity}, \emph{On Monogamy}, \emph{Prescription against Heretics},
\emph{Against Hermogenes}, \emph{Against the Valentinians}, \emph{Against
Marcion}, \emph{Against Praxeas}, \emph{On the Flesh of Christ}, \emph{On the
Soul}, \emph{On the Resurrection of the Flesh}, \emph{On Fasting}, and
\emph{On Modesty}. \emph{Against the Jews} is used with its unity and chapters
9--14 qualified. Spurious and lost works are not silently merged into this
corpus.

External controls include the Acts of the Scillitan Martyrs, \emph{Passion of
Perpetua and Felicity}, Eusebius's hostile and excerpted New Prophecy dossier in
\emph{Ecclesiastical History} 5.16--19, his other Tertullian notices,
Lactantius, Jerome, Augustine, Vincent of L\'erins, the composite Gelasian
dossier, manuscript catalogs, and the publication history of early editions.
Official Catholic reception is limited to the Catechism's actual citations and
Benedict XVI's catechesis. Modern controls include Barnes, Rankin, Osborn,
Dunn, Trevett, Wilhite, Secord, Storin and Muehlberger, Binder, current
critical editions, and manuscript and Vetus Latina research tools.

Every consequential biographical claim was tested first against Tertullian's
own corpus and the earliest reasonably accessible external witness, then
against genre, opponent, textual locus, relative chronology, transmission, and
named modern reconstruction. Self-authored evidence establishes what the work
says and how the author presents himself; it does not make polemic a neutral
transcript. A late source may preserve a lost report, but its parts are
disaggregated rather than accepted as a package. Official reception establishes
what the authority selected, not a total verdict on the author.

The negative search specifically sought contemporary biography; a birth,
baptism, ordination, excommunication, death, burial, relic, feast, or cult
record; a named Roman bishop's act; a legal career record; secure identity with
the \emph{Digest} jurist; and a controlled editorial notice for the Perpetua
passion. No such record was located. This is a bounded result from the
reasonably accessible corpus, ancient reception, catalogs, and current
scholarship, not proof that no lost record ever existed.

Internet-accessible sources and mutable institutional records were last checked
through 2026-07-19. Manuscript counts, catalog descriptions, critical editions,
and current scholarly judgments remain subject to revision.

\subsection{Limits, rights, and review state}

This release completed internal historical, theological, chronological,
transmission, tradition, and Catholic-reception audits at publication scale.
Independent patristic, textual-critical, theological, rights, and editorial
review remains outstanding. The internal work did not collate every manuscript,
verify every Latin biblical citation against the complete Vetus Latina
apparatus, inspect every modern copyrighted edition in full, reconstruct all
opponents from independent sources, or conduct an exhaustive archaeological
and epigraphic survey of Roman Carthage. No conclusion requiring those reviews
is represented as established.

Ancient Latin and Greek are public domain. Public-domain ANF and NPNF research
translations are cited and mostly paraphrased, with brief wording tied to exact
loci. Modern translations, critical editions, scholarship, Vatican English,
catalog metadata, and manuscript images retain their respective rights; online
access does not imply permission to reproduce. No third-party image appears.
This publication's original prose falls under the repository license, while
third-party source material remains outside it. The document-wide compact
reuse notice appears as the final-page colophon after generation metadata.
